\subsection{Scheduling}
\label{sec:algo_other}

In many situations, such as with causal masking or variable sequence length (varlen), the attention kernel is naturally load-imbalanced -- SMs are assigned worktiles whose mainloops differ in length, since some worktiles require more loads and MMAs than others. Furthermore, we can choose the order in which SMs process tiles, for instance by defining a preferred linearization of the grid coordinates. Abstracting away any specific features of attention, we can then apply general results on makespan minimization for identical parallel processors to our context. In particular, in FlashAttention-4, we use the classical idea of longest-processing-time-first (LPT) scheduling~\citep{graham1969}. We emphasize that the way in which we apply this idea works across all GPU architectures and was also validated as an improvement to FlashAttention-3 on Hopper GPUs.

\textbf{LPT for causal masking.} The standard attention grid is given by (mblocks, heads, batches) and is computed in increasing order left-to-right. But scores are masked out above the diagonal, so for fixed head and batch, SMs end up inefficiently processing worktiles from shortest to longest. On the other hand, a naive LPT order is also suboptimal, since for different batches, mainloop KV loads won't hit in L2 cache, and loading all KV heads first can thrash the L2 cache if they exceed its capacity. Instead, we always process batches as the outermost dimension, and swizzle over heads. This means that we divide heads into sections that don't overflow L2 cache; the tile scheduler then traverses the grid by heads per section, mblocks in reverse order, sections, and finally batches. In particular, for MQA or GQA, we always traverse all query heads per KV head before varying over mblocks. Empirically, we validate that this LPT order is highly effective; for example, for BF16 and head dimension 128 we obtain 4-8\% FLOPS gain for MHA and 7-14\% for MQA 8 as measured on an H200 GPU.

\textbf{LPT for variable sequence length.} For varlen, we also have to contend with load imbalance due to variation among batches. For instance, in decode workloads, different batches might attend to different amounts of context, and in mixed or continuous batching, some batches might be prefill while others are decode. The list of query and KV sequence lengths per batch is typically stored as attention metadata on device, and the standard varlen attention kernel reads these integers at runtime while processing batches in increasing order. However, the given batch order may be arbitrarily suboptimal with respect to load balancing -- for example, we could have shorter square prefills followed by long-context decodes. To ameliorate this, we can enforce an LPT order by launching a preprocessing kernel to sort batches according to their maximum per-worktile execution time, writing out the additional metadata of a virtual to actual batch index mapping that will be subsequently read back into the attention kernel in order to traverse batches in sorted order. This metadata can be cached and thus results in no performance loss from sorting.


